% ===========================================================================
%  Capítulo — Fundamentos del dominio y estado del arte
% ===========================================================================
\chapter{Fundamentos del dominio y estado del arte}
\label{cap:estadoarte}

La asignación de horarios es un área de estudio dentro de la planificación de
recursos y la optimización que abarca múltiples contextos prácticos~\cite{babaei2015}.
Este capítulo establece las bases conceptuales del dominio y analiza el panorama
de soluciones existentes.
La sección~\ref{sec:tipos-horarios} define los tipos de problemas de asignación
de horarios y sus variantes representativas.
La sección~\ref{sec:estadoarte-horarios} presenta el estado del arte en ambos
dominios.
Las secciones~\ref{sec:apps-existentes} y~\ref{sec:tecnologias-existentes}
examinan las herramientas comerciales disponibles y las tecnologías sobre las
que se apoyan.
Finalmente, la sección~\ref{sec:limitaciones-arte} identifica las limitaciones
que motivan el trabajo presentado en esta tesis.

% ---------------------------------------------------------------------------
\section{Problemas de asignaci\'on de horarios}
\label{sec:tipos-horarios}
% ---------------------------------------------------------------------------

La asignación de horarios constituye un problema de planificación presente en
muy diversos contextos~\cite{babaei2015}.  Consiste en organizar y distribuir
un conjunto de actividades en determinados intervalos de tiempo, bajo un
conjunto de restricciones que limitan las posibles asignaciones.

Estas restricciones se clasifican en dos tipos.  Las \textbf{restricciones
duras} son obligatorias: si no se cumplen, la solución no es válida ---dos
clases no pueden compartir el mismo aula a la misma hora; un profesor no puede
impartir dos asignaturas simultáneamente---.  Las \textbf{restricciones
blandas} no invalidan la solución pero afectan su calidad: concentrar todas
las clases de una asignatura en un solo día o colocar una materia de alta
carga cognitiva en el último turno de la tarde son situaciones admisibles
pero no deseables.

Entre las restricciones más comunes se distinguen los \textbf{conflictos} y
los \textbf{solapamientos}.  Un conflicto ocurre cuando dos o más actividades
compiten por el uso simultáneo de un mismo recurso limitado ---un aula, un
profesor, un canal de emisión---.  Un solapamiento se produce cuando
actividades asignadas coinciden en el tiempo de manera incompatible,
impidiendo su ejecución simultánea por parte de los actores involucrados.

Además de estas restricciones fundamentales, es necesario gestionar la
disponibilidad real de espacios y personas, lo que incrementa la complejidad
del problema y restringe aún más el conjunto de soluciones viables.  El
problema varía según el contexto: cada entorno tiene sus propias reglas y
objetivos.  Las subsecciones siguientes describen las variantes más relevantes
para este trabajo.

% ............................................................................
\subsection{Horario docente universitario}
\label{ssec:horario-facultad}
% ............................................................................

Este tipo de problema consiste en asignar cursos, profesores y aulas a
determinados períodos de tiempo cumpliendo múltiples restricciones~\cite{abdennadher2000,muller2005,babaei2015}:
la disponibilidad de los docentes, la capacidad de las aulas y la no
superposición de materias para un mismo grupo de estudiantes.  El objetivo
principal es generar un horario factible que optimice el uso de los recursos
y minimice conflictos.

Por ejemplo, al programar la asignatura \textit{Álgebra} para el grupo C111
los lunes a las 9:15~h en el aula~5, el sistema debe asegurar que el profesor
asignado no tenga otra clase a esa misma hora, que el aula~5 no esté ocupada
por otro grupo, y que el grupo C111 no tenga otra asignatura en ese turno.
Las restricciones concretas varían entre instituciones: en algunas la escasez
de aulas es la restricción dominante; en otras, los laboratorios especializados
crean los cuellos de botella más severos.

% ............................................................................
\subsection{Planificaci\'on de defensas de tesis}
\label{ssec:horario-defensas}
% ............................................................................

La planificación de defensas de tesis es una variante del problema de
asignación de horarios con características propias.  Cada defensa exige la
presencia simultánea de varios miembros de tribunal: tutor, oponente,
presidente, vocal y secretario.  La restricción dura crítica es la ausencia
de conflictos: un mismo miembro del tribunal no puede estar en dos defensas
programadas en el mismo slot horario.

A diferencia del horario docente, en este caso el recurso escaso no son las
aulas sino los profesores con el perfil adecuado para cada rol~\cite{nardo2019}.
Detectar estos conflictos de forma visual es precisamente uno de los
requisitos que motivaron el desarrollo del DSL descrito en este trabajo.

% ............................................................................
\subsection{Programaci\'on televisiva}
\label{ssec:horario-tv}
% ............................................................................

En la confección de la parrilla televisiva, el reto es organizar los programas
en franjas horarias para maximizar la audiencia y cumplir los compromisos
comerciales~\cite{bollapragada2002}.  Se consideran factores como el tipo de
público en cada franja, la competencia con otros canales y la duración exacta
de los programas~\cite{fernandez2008}.

El programador debe, por ejemplo, evitar que un noticiero importante coincida
con la transmisión en vivo de un evento deportivo de duración variable, o
garantizar que un bloque infantil se emita en las mañanas de fin de semana y
no en la franja nocturna.  A diferencia del entorno académico, el enfoque
aquí es estratégico y comercial: las restricciones blandas ---cómo posicionar
los programas para maximizar la cuota de pantalla--- tienen tanto peso como
las restricciones duras.

% ---------------------------------------------------------------------------
\section{Estado del arte}
\label{sec:estadoarte-horarios}
% ---------------------------------------------------------------------------

% ............................................................................
\subsection{Horario docente universitario}
\label{ssec:arte-facultad}
% ............................................................................

La investigación sobre generación automática de horarios universitarios tiene
décadas de historia.  Abdennadher y Marte~\cite{abdennadher2000} formalizaron
el problema mediante Reglas de Manejo de Restricciones (CHR), estableciendo
un lenguaje declarativo de restricciones que influyó en trabajos posteriores.
Müller~\cite{muller2005} desarrolló un sistema completo de timetabling basado
en satisfacción de restricciones para entornos universitarios.  Reddy
et al.~\cite{reddy2019} propusieron un modelo de coloración de grafos que
garantiza formalmente la ausencia de conflictos pero resulta computacionalmente
costoso en instancias reales.

En la Facultad de Matemática y Computación de la Universidad de La Habana
(MATCOM/UH) existe una línea de investigación consolidada sobre este problema.
Portela Aguilera~\cite{portela2010} realizó la primera aproximación formal a
la generación automática de horarios para la facultad, sentando las bases del
problema en el contexto local.  Travieso Sosa~\cite{travieso2012} formalizó
el conjunto de restricciones del problema e introdujo el concepto de
``felicidad'' como medida ponderada de la calidad del horario.  Chang
Alcover~\cite{chang2014} aplicó satisfacción de restricciones (CSP) con
retroceso y búsqueda local.  Fernández Llanos~\cite{fernandez_llanos2014}
desarrolló una API especializada con evaluación incremental de restricciones.

Seguí Bonilla~\cite{segui2014} y Tomás López~\cite{tomaslopez2014}
desarrollaron módulos de búsqueda automática de horarios basados en técnicas
de optimización computacional.  Quintero Rojas~\cite{quintero2014} integró
los módulos anteriores en un sistema completo de gestión de horarios.  Nardo
González~\cite{nardo2019} extendió la investigación al problema de
planificación de defensas de tesis, que comparte la estructura de conflictos
del horario docente pero con restricciones específicas sobre los miembros de
tribunal.  Casteleiro Wong~\cite{casteleiro2022} abordó la informatización
de la gestión docente en el contexto de la facultad.

% ............................................................................
\subsection{Programaci\'on televisiva}
\label{ssec:arte-tv}
% ............................................................................

En el ámbito de la radiodifusión, Bollapragada et al.~\cite{bollapragada2002}
plantearon el problema de la programación televisiva como programación lineal
entera mixta (MILP) orientada a maximizar la audiencia en cadenas de
televisión comerciales.  Fernández y Salido~\cite{fernandez2008} lo modelaron
como un CSP y propusieron técnicas de descomposición del espacio de búsqueda
junto con métodos de búsqueda automática.

A partir de 2015 la incorporación de técnicas de aprendizaje automático marcó
una nueva dirección.  Kim et al.~\cite{kim2021} propusieron un sistema de
aprendizaje por refuerzo profundo para la programación televisiva, publicado
en \textit{IEEE Transactions on Broadcasting}.

En ambos dominios la tendencia reciente es hacia herramientas interactivas
que asisten al experto sin reemplazarlo: el planificador mantiene el control
creativo mientras el sistema garantiza la viabilidad técnica frente a las
restricciones.

% ---------------------------------------------------------------------------
\section{Aplicaciones existentes para la gesti\'on de horarios}
\label{sec:apps-existentes}
% ---------------------------------------------------------------------------

Dada la complejidad y la gran cantidad de variables de estos problemas,
confeccionar horarios de forma manual en herramientas básicas resulta un
proceso lento y propenso a errores humanos~\cite{pillay2016}.  Por este
motivo existen aplicaciones especializadas, tanto en el ámbito académico
como en el de la radiodifusión.  Estas aplicaciones van desde soluciones de
escritorio hasta plataformas de software como servicio (\textit{Software as
a Service}, SaaS) accesibles desde el navegador.

% ............................................................................
\subsection{\'Ambito acad\'emico}
\label{ssec:apps-academico}
% ............................................................................

En la gestión del \textit{Educational Timetabling Problem}
(ETP)~\cite{babaei2015} destacan las siguientes herramientas:

\begin{description}
  \item[Coursedog y TimeEdit.] Representan el estado del arte en educación
    superior~\cite{bass2021}.  Son plataformas SaaS que permiten a
    coordinadores y decanos gestionar recursos en tiempo real, con interfaces
    de validación que alertan visualmente sobre conflictos de aula o
    solapamientos de profesores.

  \item[aSc Horarios.] Una de las herramientas más extendidas para educación
    primaria y secundaria~\cite{pillay2016}.  Destaca por su capacidad de
    generar horarios automáticos basados en heurísticas y por su interfaz
    interactiva de arrastrar y soltar que facilita la edición manual asistida.

  \item[Syllabus Plus.] Utilizada por grandes universidades, se centra en la
    optimización del horario desde la perspectiva del estudiante: garantiza
    que el calendario sea viable para el cuerpo estudiantil y no solo para
    el cuerpo docente.

  \item[Horarium.ai.] Plataforma de generación automática e interactiva de
    horarios mediante inteligencia artificial~\cite{horarium2024}.  Incorpora
    modelos de lenguaje para asistir al planificador en la toma de decisiones
    durante la confección del horario.
\end{description}

% ............................................................................
\subsection{\'Ambito de radiodifusi\'on}
\label{ssec:apps-tv}
% ............................................................................

En televisión y radiodifusión, las aplicaciones deben manejar la precisión
temporal de la emisión junto con las restricciones comerciales:

\begin{description}
  \item[MEDIAGENIX (WHATS'On).] Estándar europeo para la gestión de canales
    de televisión y plataformas VOD.  Permite la planificación de parrillas a
    largo plazo, validando automáticamente los derechos de emisión y las
    reglas de pauta publicitaria~\cite{mediagenix2024}.

  \item[BroadView y Amagi.] Soluciones que integran la planificación del
    tráfico de emisión con la automatización del \textit{playout}.  Aseguran
    que la lista de programación se ejecute sin errores técnicos ni huecos
    de emisión (\textit{blackouts})~\cite{matthews2018}.
\end{description}

% ---------------------------------------------------------------------------
\section{Tecnolog\'{\i}as empleadas}
\label{sec:tecnologias-existentes}
% ---------------------------------------------------------------------------

La arquitectura tecnológica de estas herramientas ha evolucionado hacia
ecosistemas en capas que separan la interfaz de usuario, el motor de
optimización y el almacenamiento:

\begin{description}
  \item[Interfaces de usuario.] La mayoría de las aplicaciones modernas
    utilizan \textit{frameworks} de JavaScript como React.js, Vue.js o
    Angular~\cite{garcia2024}.  Estas tecnologías permiten crear interfaces
    reactivas donde el usuario recibe retroalimentación inmediata ---cambios
    de color, alertas--- ante el incumplimiento de una restricción.

  \item[Backend y motores de optimizaci\'on.] El procesamiento intensivo se
    delega a servicios escritos en lenguajes de alto rendimiento como Java,
    Go o Python.  Estos servicios ejecutan motores de satisfacción de
    restricciones (\textit{Constraint Solvers}) como OptaPlanner o Google
    OR-Tools, que operan de forma independiente a la interfaz de usuario.

  \item[Persistencia.] Las aplicaciones utilizan bases de datos relacionales
    como PostgreSQL o SQL Server para gestionar la concurrencia cuando varios
    usuarios editan el horario simultáneamente y para mantener un historial
    de auditoría de los cambios.

  \item[Formatos de intercambio.] Para la exportación de datos, los sistemas
    académicos emplean estándares como JSON e iCal (.ics).  En el ámbito
    televisivo, el estándar predominante es el BXF (\textit{Broadcast
    eXchange Format}), un esquema XML que permite la comunicación entre el
    software de programación y los equipos de emisión física~\cite{smpte2021}.
    En entornos institucionales más modestos, el formato de intercambio más
    frecuente es la hoja de cálculo: un archivo \xlsx{} que concentra toda
    la información del horario y que el personal puede editar sin software
    especializado.
\end{description}

% ---------------------------------------------------------------------------
\section{Limitaciones de los enfoques actuales}
\label{sec:limitaciones-arte}
% ---------------------------------------------------------------------------

Los sistemas descritos resuelven bien el problema de \emph{asignación}: dado
un conjunto de recursos, actividades y restricciones, encontrar una
distribución factible o cercana al óptimo.  Sin embargo, presentan una
limitación estructural cuando el requisito es producir una
\emph{visualización} del horario resultante: la representación visual queda
acoplada al sistema que la generó.

En entornos donde la herramienta de visualización preferida es la hoja de
cálculo ---porque es accesible sin instalación, editable por el usuario y
fácilmente compartible--- la generación automática de ese documento exige
código que entremezcla la lógica del problema con los detalles de la API de
la biblioteca de salida.  Esta mezcla, descrita en detalle en la
introducción, es la motivación central del lenguaje de dominio específico
desarrollado en este trabajo.
